% ATTEMPT_STATUS: unresolved
% TOP_PROBLEM: {"problemId": "problem.two-dimensional-jacobian-conjecture-in-characteristic-zero", "canonicalTitle": "Two-Dimensional Jacobian Conjecture in Characteristic Zero", "releaseRank": 103, "primaryDomain": "algebra_representation_category", "primaryDomainLabel": "Algebra, representation & category theory", "displayStatus": "open", "releaseStatus": "open"}
% !TeX program = lualatex
% Complete problem section copied verbatim from research_batch_101_103.tex.
\documentclass[11pt]{article}
\usepackage[margin=25mm]{geometry}
\usepackage{fontspec}
\setmainfont{Latin Modern Roman}
\usepackage{amsmath,amssymb,mathtools}
\usepackage{unicode-math}
\setmathfont{Latin Modern Math}
\usepackage{enumitem,needspace}
\usepackage{xurl,hyperref}
\hypersetup{colorlinks=true,urlcolor=blue}
\setcounter{secnumdepth}{1}
\setlength{\parindent}{0pt}
\setlength{\parskip}{5pt}
\setlength{\emergencystretch}{3em}
\setlist[itemize]{leftmargin=3em,itemsep=5pt,topsep=4pt,parsep=0pt}
\allowdisplaybreaks
\newcommand{\N}{\mathbb{N}}
\newcommand{\Z}{\mathbb{Z}}
\newcommand{\Q}{\mathbb{Q}}
\newcommand{\R}{\mathbb{R}}
\newcommand{\C}{\mathbb{C}}
\newcommand{\F}{\mathbb{F}}
\newcommand{\A}{\mathbb{A}}
\newcommand{\PP}{\mathbb P}
\newcommand{\E}{\mathbb{E}}
\newcommand{\eps}{\varepsilon}
\newcommand{\dd}{\,\mathrm d}
\newcommand{\sref}[2]{\hyperlink{source:#1:#2}{[S#2]}}
\newcommand{\eref}[2]{\hyperlink{extra:#1:#2}{[E#2]}}
% Turn the next switch off for a shorter reading copy. The quotations preserve
% every exactTarget field, including qualifications omitted from a short summary.
\newif\ifshowcatalogue
\showcataloguetrue
\newcommand{\cataloguescope}[1]{\ifshowcatalogue
 \paragraph{Exact scope in the supplied catalogue.}
 \begin{quote}\small #1\end{quote}\fi}
\begin{document}
\setcounter{section}{102}
% ================================================================
% problemId: problem.two-dimensional-jacobian-conjecture-in-characteristic-zero
% source releaseRank: 103; source releaseStatus: open
% exactTarget SHA-256: 8ea39d9d3dc7447f107db63684a8c632dfb14985b3201d5472b02cc82b029d57
\clearpage
\section{\texorpdfstring{Two-Dimensional Jacobian Conjecture in Characteristic Zero}{Two-Dimensional Jacobian Conjecture in Characteristic Zero}}\label{103}
\subsection{Definitions and mathematical statement}
For $F=(f,g)$ with $f,g\in k[x,y]$, its Jacobian determinant is
\[
 JF=\frac{\partial f}{\partial x}\frac{\partial g}{\partial y}
       -\frac{\partial f}{\partial y}\frac{\partial g}{\partial x}.
\]
A Keller map satisfies $JF\in k^\times$. In characteristic zero, the question asks whether there are $u,v\in k[x,y]$ making $G=(u,v)$ a two-sided polynomial inverse: $G\circ F=F\circ G=\operatorname{id}$. This is the affine plane case over every characteristic-zero field. Statements or alleged counterexamples in more variables do not decide this selected two-variable target.
\par\smallskip\textit{Formulation and concept references:} \sref{103}{1}.
\cataloguescope{Let k be a field of characteristic zero and let F:k\textasciicircum{}2 -> k\textasciicircum{}2 be a polynomial map whose Jacobian determinant is a nonzero constant. Determine whether F must have a polynomial inverse, equivalently whether every two-variable Keller map is a polynomial automorphism of the affine plane.}
\subsection{Short English statement}
Does every plane polynomial map with constant nonzero Jacobian determinant have a polynomial inverse?
% BEGIN RESEARCH ADDITION ATTEMPT-103
\subsection{Research attempt}

\paragraph{Current frontier.}
The selected statement is the plane problem over every characteristic-zero
field. Theorem 2.1 in the degree analysis of Guccione, Guccione, Horruitiner,
and Valqui excludes a plane counterexample of total degree below $108$.
This is a finite-degree result, not an argument for arbitrary degrees.
\eref{103}{1}
The 2026 sources \sref{103}{1}--\sref{103}{3} concern constructions in
higher dimension and explicitly distinguish the plane case. Rather than
accepting a solution claim merely from its title, the particular
three-variable formula used below was checked by exact polynomial expansion.
That check is separate from, and does not certify, the other results of the
cited manuscripts.

\paragraph{Attack route.}
A properness argument would have to exclude branches escaping to infinity
with bounded image; a nonzero Jacobian supplies local invertibility, not that
global control. A formal-inverse argument would have to prove termination of
the inverse series, not merely its existence. A counterexample route is to
restrict the higher-dimensional example to a plane and retain two output
coordinates. We test that route first, then pursue a different one:
leading-coefficient identities, rational normal forms, and the polynomiality
constraints at their apparent poles. All new algebra below works over an
arbitrary field $k$ of characteristic zero.

\paragraph{Attempt: testing a higher-dimensional construction.}
The following map is the example displayed in equation (1) of
\sref{103}{1}, also discussed in Section 3 of \sref{103}{2}. Write $a=1+xy$:
\begin{equation}\label{ra103:threeD}
 \begin{aligned}
 F_1&=a^3z+y^2a(4+3xy),\\
 F_2&=y+3xa^2z+3xy^2(4+3xy),\\
 F_3&=2x-3x^2y-x^3z.
 \end{aligned}
\end{equation}
Exact expansion in $\Q[x,y,z]$ and rational substitution give
\begin{align*}
 \det DF&=-2,\\
 F(0,0,-1/4)&=F(1,-3/2,13/2)=F(-1,3/2,13/2)
              =(-1/4,0,0).
\end{align*}
These identities verify a noninjective Keller map in \emph{three} variables;
they are not a counterexample to the statement of this section.

The last two preimages lie on the plane $z=13/2$. Restricting to that plane
and retaining any two output coordinates preserves their collision, so this
is a genuine candidate-generating operation. It fails the Jacobian
hypothesis. If $J_{ij}$ denotes the two-variable Jacobian of the restricted
pair $(F_i,F_j)$, exact evaluation gives
\begin{center}
\begin{tabular}{c|r r}
$(i,j)$ & $J_{ij}(0,0)$ & $J_{ij}(1,0)$\\
\hline
$(1,2)$ & $0$ & $-1521/4$\\
$(1,3)$ & $0$ & $1365/4$\\
$(2,3)$ & $-2$ & $1283/2$
\end{tabular}
\end{center}
Thus none is a plane Keller pair. This calculation tests a specified
restriction, not every possible nonlinear reduction in dimension.

A direct plane tangent sweep is also obstructed. For a polynomial curve
$C(u)\in k[u]^2$, put $S(u,t)=C(u)+tC'(u)$. Then
\[
 \det\left(\frac{\partial S}{\partial u},
           \frac{\partial S}{\partial t}\right)
     =t\det(C''(u),C'(u)),
\]
which vanishes on $t=0$. Polynomial automorphisms before or after this map
multiply its Jacobian by nonzero constants and cannot remove this critical
set. A singular rational coordinate change could change the situation, but
its denominators and exceptional loci would then require a separate proof.
This motivates checking polynomiality, rather than only rational Jacobian
identities, in the next approach.

\paragraph{Leading-degree reduction.}
Write the plane variables as $u,t$, and suppose
\[
 J(f,g):=f_u g_t-f_t g_u=\eta\in k^\times.
\]
If $r=\deg_t f>0$ and $s=\deg_t g>0$, write their leading coefficients as
$A(u)$ and $B(u)$. The coefficient of $t^{r+s-1}$ in the Jacobian must vanish:
\begin{equation}\label{ra103:leading}
 sA'B-rAB'=0.
\end{equation}
In particular, if $r\mid s$, then
$(B/A^{s/r})'=0$ in $k(u)$, hence $B/A^{s/r}=c\in k^\times$.
The target shear
\begin{equation}\label{ra103:shear}
 (f,g)\longmapsto(f,g-cf^{s/r})
\end{equation}
strictly reduces $\deg_t g$ and preserves the Jacobian. It is an invertible
polynomial change of target coordinates. The symmetric operation applies
when $s\mid r$.

If one $t$-degree is zero, say $f=f(u)$, then $f'(u)g_t=\eta$. The only
units of $k[u,t]$ are nonzero constants, so $f'(u)$ and $g_t$ are constants.
Characteristic zero gives
$f=\alpha u+\beta$ and $g=\gamma t+h(u)$ with
$\alpha,\gamma\in k^\times$. The inverse is explicitly
\[
 u=\frac{f-\beta}{\alpha},\qquad
 t=\frac{g-h((f-\beta)/\alpha)}{\gamma}.
\]
Thus the divisible-degree route terminates in an automorphism whenever it
can be continued. What remains is to handle nondividing degree pairs.

\paragraph{A worked obstruction: the pair of $t$-degrees $(2,3)$.}
Suppose
\[
 f=a(u)t^2+b(u)t+c_0(u),\qquad
 g=d(u)t^3+e(u)t^2+h(u)t+j(u),
\]
with $a,d\ne0$. Equation \eqref{ra103:leading} says $3a'd-2ad'=0$,
so $d^2/a^3\in k^\times$. Unique factorization in $k[u]$ implies
\[
 a=Aq^2,\qquad d=Dq^3
 \quad\text{for }q\in k[u]\setminus\{0\},\ A,D\in k^\times.
\]
Indeed, for each irreducible factor the valuations satisfy
$2\operatorname{ord}(d)=3\operatorname{ord}(a)$; those of $a,d$ are
respectively twice and three times a nonnegative integer. Constants need
not have square or cube roots in $k$.

Scale the two output coordinates by $A^{-1},D^{-1}$, relabel coefficients,
and retain the symbol $\eta$ for the resulting nonzero constant Jacobian.
Set
\[
 v=qt+r_0,\qquad r_0=\frac{b}{2q}\in k(u),\qquad
 P=c_0-r_0^2.
\]
This is an auxiliary change in $k(u)[t]$, \emph{not yet a polynomial
coordinate change}. In these coordinates
\[
 f=v^2+P(u),\qquad
 g=v^3+L(u)v^2+E(u)v+H(u),
\]
where $L,E,H$ are rational functions of $u$. The chain rule yields
$J_{u,t}(f,g)=qJ_{u,v}(f,g)$, and direct differentiation gives
\begin{equation}\label{ra103:cubiccompare}
 J_{u,v}(f,g)=
 -2L'v^3+(3P'-2E')v^2+(2LP'-2H')v+P'E.
\end{equation}
Since this must equal $\eta/q$, coefficient comparison gives
\[
 L=\ell\in k,\qquad E=\frac32P+\kappa,\qquad
 H=\ell P+\mu\quad(\kappa,\mu\in k).
\]
Subtract the polynomial target expression $\ell f+\mu$ from $g$. The new
coordinate $g_*$ is still in $k[u,t]$ and satisfies
\begin{equation}\label{ra103:cubicnormal}
 g_*=v^3+\left(\frac32P+\kappa\right)v,
 \qquad qP'\left(\frac32P+\kappa\right)=\eta.
\end{equation}

Now the integrality constraint becomes decisive. If $r_0$ had a pole of
order $m>0$ at an irreducible polynomial of $k[u]$, then
$P=c_0-r_0^2$, and the constant coefficient in $t$ of $g_*$ would be
\begin{equation}\label{ra103:pole}
 r_0^3+\left(\frac32(c_0-r_0^2)+\kappa\right)r_0
 =-\frac12r_0^3+\left(\frac32c_0+\kappa\right)r_0.
\end{equation}
The first term has pole order $3m$; the second has pole order at most $m$.
They cannot cancel. This contradicts $g_*\in k[u,t]$. Therefore $r_0$ has
no finite poles, so $r_0\in k[u]$, and then $P\in k[u]$ as well.

If $P$ is constant, the last equation of \eqref{ra103:cubicnormal} is zero
on the left. If $\deg P=e\ge1$, its left side has degree
$\deg q+2e-1\ge1$. Neither possibility equals the nonzero constant $\eta$.
We have proved that a Keller pair cannot have $t$-degrees $(2,3)$.
\hfill$\square$

\paragraph{Pushing the obstruction through every odd partner degree.}
The previous calculation extends beyond the cubic case. This extension is
proved here to avoid relying on a pattern inferred from small computations.
Suppose $\deg_t f=2$ and $\deg_t g=2m+1$ with $m\ge1$. The leading identity
now gives, after constant rescaling,
\[
 f=v^2+P(u),\qquad v=q(u)t+r_0(u),\qquad
 g=v^{2m+1}+\text{lower powers of }v,
\]
with $q\in k[u]\setminus\{0\}$, $r_0,P\in k(u)$, and
$P=c_0-r_0^2$ for $c_0\in k[u]$. Define, for $\ell\ge0$,
\begin{equation}\label{ra103:oddforms}
 H_\ell(v,P)=\sum_{j=0}^{\ell}
   \binom{\ell+1/2}{j}P^jv^{2\ell+1-2j},
 \qquad
 c_\ell=\binom{\ell+1/2}{\ell}\ne0.
\end{equation}
The generalized binomial coefficients here are rational numbers, viewed in
$k$. For the derivation
$\mathcal D=P'\partial_v-2v\partial_u$, where $\partial_u$ holds $v$ fixed,
coefficient cancellation gives
\begin{equation}\label{ra103:derivation}
 \mathcal D(H_\ell)=c_\ell P'P^\ell,
 \qquad \mathcal D((v^2+P)^j)=0.
\end{equation}
To verify the first identity in all degrees, the coefficient of each positive
power of $v$ cancels using
$(j+1)\binom{\ell+1/2}{j+1}
 =(\ell+1/2-j)\binom{\ell+1/2}{j}$; only the derivative of the final
$v$ term contributes a constant. Thus no induction is being inferred solely
from the first few cases.

Any polynomial $G\in k(u)[v]$ with $\mathcal D G\in k(u)$ has the form
\begin{equation}\label{ra103:classification}
 G=\sum_{\ell}\alpha_\ell H_\ell(v,P)
      +\sum_j\beta_j(v^2+P)^j,
 \qquad\alpha_\ell,\beta_j\in k.
\end{equation}
Here is a degree induction proving it. If $a(u)v^n$ is the leading term,
the coefficient of $v^{n+1}$ in $\mathcal DG$ is $-2a'(u)$, so $a\in k$.
For even $n$, subtract $a(v^2+P)^{n/2}$; for odd $n=2\ell+1$, subtract
$aH_\ell$. In each case the degree decreases and the derivative of the
remainder is still constant in $v$. The degree-zero remainder also has
constant coefficient, since its derivative contributes $-2a'v$.
This completes the induction, using the characteristic-zero fact that the
constants of differentiation on $k(u)$ are exactly $k$.

Apply this to $g$, for which $\mathcal Dg=\eta/q$. Its leading coefficient
is $1$, so $\alpha_m=1$. Remove the even part in
\eqref{ra103:classification} by subtracting a polynomial in $f$ with
coefficients in $k$. The resulting $g_*$ remains a polynomial in $u,t$ and
is $\sum_{\ell=0}^m\alpha_\ell H_\ell$.
If $r_0$ has a pole of order $b>0$, then the highest pole in its constant
coefficient in $t$ is supplied by $H_m(r_0,c_0-r_0^2)$. Its leading
coefficient is
\begin{equation}\label{ra103:binomialpole}
 \sum_{j=0}^m(-1)^j\binom{m+1/2}{j}
       =(-1)^m\binom{m-1/2}{m}\ne0.
\end{equation}
The equality follows by telescoping Pascal's identity. Terms involving a
positive power of the regular function $c_0$ have smaller pole order, and
$H_\ell$ with $\ell<m$ has order at most $(2\ell+1)b<(2m+1)b$.
The leading pole cannot cancel. Again $r_0,P\in k[u]$.

Finally \eqref{ra103:derivation} gives
\begin{equation}\label{ra103:oddcontradiction}
 qP'R(P)=\eta,
 \qquad R(Z)=\sum_{\ell=0}^m\alpha_\ell c_\ell Z^\ell\in k[Z],
 \quad \deg R=m.
\end{equation}
For constant $P$ the left side vanishes. For $\deg P=e\ge1$ it has degree
$\deg q+(m+1)e-1\ge1$. This proves that no Keller pair has $t$-degrees
$(2,2m+1)$ for any $m\ge1$.

\textbf{Restricted invertibility lemma.}
If one component of a plane Keller map over a characteristic-zero field
has degree at most two in one fixed source variable, then the map is a
polynomial automorphism. In particular, this holds if both components have
degree at most three in that variable. There is no restriction on their
degrees in the other variable.

\emph{Proof.} Degree zero was handled above. With a component of degree one,
repeated use of \eqref{ra103:shear} decreases the partner degree until the
degree-zero case applies. With a component of degree two, an even partner
degree is reduced by the same shear. An odd partner degree greater than one
has just been ruled out; degree one or zero leads to the previous cases.
For two degree-three components, one constant linear combination cancels
the leading coefficient by \eqref{ra103:leading}, leaving a component of
degree at most two. Every operation used is a polynomial automorphism of the
target, so reversing the finite sequence supplies a two-sided polynomial
inverse over the original field $k$.\hfill$\square$

This is a restricted theorem proved within the notebook, not a claim of a
new result in the literature. It is genuinely not a total-degree-three
argument: for any $M\ge1$,
\[
 f=t+u^M,\qquad g=(t+u^M)^3+u
\]
has $J(f,g)=-1$, $t$-degrees $(1,3)$, and inverse
$u=g-f^3$, $t=f-(g-f^3)^M$. Its total degree grows with $M$.

\paragraph{The next unresolved degree pair: an explicit system.}
To test whether the method really generalizes, consider $t$-degrees $(3,4)$.
The leading relation permits normalizing in $k(u)[t]$ to
\[
 f=v^3+P(u)v+Q(u),\qquad v=q(u)t+r_0(u),\quad q\in k[u]\setminus\{0\}.
\]
After removing a constant multiple of $f$ from $g$ to cancel its $v^3$ term,
coefficient comparison forces
\begin{align}\label{ra103:quartic}
 g={}&v^4+\left(\frac43P+a\right)v^2
           +\left(\frac43Q+b\right)v\\
    &+\frac29P^2+\frac23aP+c,\notag
 \qquad a,b,c\in k,
\end{align}
provided $J_{u,v}(f,g)$ is constant in $v$. The remaining two conditions are
\begin{align}\label{ra103:nextsystem}
 \frac43PQ+bP+2aQ&=d_0\in k,\\
 q\left[\left(\frac43Q+b\right)Q'
      -\left(\frac49P^2+\frac23aP\right)P'\right]&=\eta.\notag
\end{align}
For example, the $v$ coefficient of the Jacobian is the derivative of the
first expression, and its constant coefficient is the bracketed expression
in the second. The higher coefficients cancel. These assertions were also
checked by exact symbolic expansion in the accompanying script.

The first equation can be rewritten
\[
 \left(P+\frac32a\right)\left(Q+\frac34b\right)
       =\frac34d_0+\frac98ab.
\]
Thus there is extra rigidity, but it is a rational-function constraint, not
a proof of polynomiality. Both original coordinates must still lie in
$k[u,t]$ after $v=qt+r_0$ is substituted. Here the potential poles of
$r_0,P,Q$ interact; the single uncancellable term from the quadratic case
has not been proved to survive. Nor have all rational solutions of
\eqref{ra103:nextsystem} compatible with those integrality conditions been
excluded. This is a concrete next obstruction, rather than an asserted
induction step.

\paragraph{Stress test.}
\emph{Field and denominator issues.}
The pole arguments use valuations at irreducible polynomials of $k[u]$,
not the existence of roots in $k$ or an embedding into $\C$. Constant
rescalings avoid adjoining roots of leading constants. The substitutions
$v=qt+r_0$ are used only over $k(u)$; they were never called polynomial
automorphisms. Invertibility of the restricted maps follows instead from
the actual target shears and the triangular terminal case.

\emph{A false global induction.}
Equation \eqref{ra103:leading} by itself does not imply that $r\mid s$ or
$s\mid r$. For example, the leading coefficients $A=q^3$, $B=q^4$
satisfy it for $(r,s)=(3,4)$. The lower-coefficient equations and integrality
are indispensable. A sufficient missing lemma would assert that every plane
Keller pair with positive $t$-degrees has a divisible degree pair, allowing
\eqref{ra103:shear} to continue. Iterating that lemma would prove the entire
plane conjecture, so treating it as a routine degree fact would be circular.
The restricted result proves only some cases of the needed rigidity.

\emph{Characteristic and computational checks.}
Characteristic zero is essential. In characteristic $p>0$, the map
$(u,t)\mapsto(u-u^p,t)$ has Jacobian $1$ but identifies $(0,0)$ and $(1,0)$.
The claims about constants of differentiation and integration of a constant
derivative would fail in that setting. The script checks
\eqref{ra103:cubiccompare}, \eqref{ra103:pole},
\eqref{ra103:derivation} for $0\le\ell\le6$,
\eqref{ra103:binomialpole} in the same range, and the coefficient system
\eqref{ra103:nextsystem}. The all-degree binomial proof above, not the finite
checks, justifies the odd-degree extension. None of the computations searches
all polynomial maps or constitutes a formal proof-assistant verification.

\paragraph{Outcome.}
\textbf{Outcome: Exploratory approach with unresolved gap.}

The counterexample-oriented plane restrictions fail the constant-Jacobian
hypothesis, with exact witnesses to that failure. On the proof side, the
leading-degree reduction and pole argument establish the restricted
invertibility lemma, including arbitrary odd partner degrees when one
component is quadratic in a fixed variable. The first unresolved $(3,4)$
normal form has been reduced to the explicit rational system
\eqref{ra103:nextsystem} together with polynomiality constraints.

No argument supplied here handles every nondividing higher-degree pair, and
no plane counterexample has been produced. A complete proof would need a
uniform replacement for the missing divisibility or pole-rigidity step; a
counterexample would have to satisfy both the nonzero constant Jacobian and
global polynomiality, not merely the rational normal-form equations. No
priority claim is made for the restricted lemma or its proof.

% END RESEARCH ADDITION ATTEMPT-103
\subsection{Sources}
\begin{itemize}\raggedright
\item[\textnormal{[S1]}]\hypertarget{source:103:1}{} Terence Tao, "A digestion of the Jacobian conjecture counterexample", What's new, 21 July 2026.
\url{https://terrytao.wordpress.com/2026/07/21/a-digestion-of-the-jacobian-conjecture-counterexample/}.
{\small\textit{Catalogue formulation source.}}
\item[\textnormal{[S2]}]\hypertarget{source:103:2}{} Shuhong Gao, Counterexamples to the Jacobian conjecture in dimensions greater than two.
\url{https://arxiv.org/abs/2608.00222}.
{\small\textit{Catalogue research/status reference; not a proof endorsement.}}
\item[\textnormal{[S3]}]\hypertarget{source:103:3}{} Gao, section 2: Background and previous work.
\url{https://arxiv.org/html/2608.00222v1#S2}.
{\small\textit{Catalogue research/status reference; not a proof endorsement.}}
% BEGIN RESEARCH ADDITION REFERENCES-103
\item[\textnormal{[E1]}]\hypertarget{extra:103:1}{}
J. A. Guccione, J. J. Guccione, R. Horruitiner, and C. Valqui,
\emph{Increasing the degree of a possible counterexample to the Jacobian
Conjecture from 100 to 108}, arXiv:2204.14178v1, 29 April 2022.
Theorem 2.1 and Sections 2--6; the total-degree exclusion is distinct from
this notebook's degree-in-one-variable calculation.
\url{https://arxiv.org/html/2204.14178v1}.
{\small\textit{Consulted 25 September 2026; cited as a partial-result
manuscript, not as a proof of the full conjecture.}}
% END RESEARCH ADDITION REFERENCES-103
\end{itemize}

\end{document}
